\documentclass[addpoints,spanish, 12pt,a4paper]{exam}
%\documentclass[answers, spanish, 12pt,a4paper]{exam}
\printanswers
\pointpoints{punto}{puntos}
\hpword{Puntos:}
\vpword{Puntos:}
\htword{Total}
\vtword{Total}
\hsword{Resultado:}
\hqword{Ejercicio:}
\vqword{Ejercicio:}

\usepackage[utf8]{inputenc}
\usepackage[spanish]{babel}
\usepackage{eurosym}
%\usepackage[spanish,es-lcroman, es-tabla, es-noshorthands]{babel}


\usepackage[margin=1in]{geometry}
\usepackage{amsmath,amssymb}
\usepackage{multicol}
\usepackage{yhmath}

\pointsinrightmargin % Para poner las puntuaciones a la derecha. Se puede cambiar. Si se comenta, sale a la izquierda.
\extrawidth{-2.4cm} %Un poquito más de margen por si ponemos textos largos.
\marginpointname{ \emph{\points}}

\usepackage{graphicx}

\graphicspath{{../../img/}} 

\newcommand{\class}{2º Bachillerato CIT}
\newcommand{\examdate}{\today}
\newcommand{\examnum}{Parcial 2ªEv.}
\newcommand{\tipo}{A}


\newcommand{\timelimit}{45 minutos}

\renewcommand{\solutiontitle}{\noindent\textbf{Solución:}\enspace}


\pagestyle{head}
\firstpageheader{\includegraphics[width=0.2\columnwidth]{header_left}}{\textbf{Departamento de Matemáticas\linebreak \class}\linebreak \examnum}{\includegraphics[width=0.1\columnwidth]{header_right}}
\runningheader{\class}{\examnum}{Página \thepage\ of \numpages}
\runningheadrule


\usepackage{pgf,tikz,pgfplots}
\pgfplotsset{compat=1.15}
\usepackage{mathrsfs}
\usetikzlibrary{arrows}


\begin{document}

\noindent
\begin{tabular*}{\textwidth}{l @{\extracolsep{\fill}} r @{\extracolsep{6pt}} }
\textbf{Nombre:} \makebox[3.5in]{\hrulefill} & \textbf{Fecha:}\makebox[1in]{\hrulefill} \\
 & \\
\textbf{Tiempo: \timelimit} & Tipo: \tipo 
\end{tabular*}
\rule[2ex]{\textwidth}{2pt}
Esta prueba tiene \numquestions\ ejercicios. La puntuación máxima es de \numpoints. 
La nota final de la prueba será la parte proporcional de la puntuación obtenida sobre la puntuación máxima. 

\begin{center}


\addpoints
 %\gradetable[h][questions]
	\pointtable[h][questions]
\end{center}

\noindent
\rule[2ex]{\textwidth}{2pt}

\begin{questions}

%\question 
%
%\begin{parts}
%\part[2] 
%\begin{solution}
%\end{solution}
%
%
%\end{parts}
%\addpoints

\question[1] Determine los valores de $a$ y $b$ para sea
continua la función:
$$f(x)=\left\{\begin{array}{ll}
\dfrac{1}{e^x} & si \ x\leq 0 \\
a\cos(x) + b & si \  0 < x \leq \pi \\
\sen(x)-ax & si \ \pi<x
\end{array}\right.$$
\begin{solution}
Fuera de $x=0$ y $x=\pi$ la función es continua por serlo sus trozos en sus dominios. \\

En $x=0$: \\
$\lim_{x \to 0^-}f=\frac{1}{e^0}=1$\\
$\lim_{x \to 0^+}f=a\cdot \cos(0) +b =a+b \to a+b=1$\\

En $x=\pi$: \\
$\lim_{x \to \pi^-}f=a\cdot \cos(\pi) +b =b-a$\\
$\lim_{x \to \pi^+}f=\sen(\pi) - a\cdot \pi =- a\cdot \pi \to b-a=-a\cdot \pi$\\

$(2-\pi)\cdot a= 1 \to a=\frac{1}{2- \pi} \to b=\frac{1-\pi}{2-\pi}$

\end{solution}

\question Calcula las siguientes integrales:
\begin{parts}
\part[1]
$$\int x^{2} \ln{\left(x \right)}\, dx$$
\begin{solution}
$\left|\begin{matrix}
u = \ln x & du=\frac{1}{x} \, dx \\ dv= x^2 \, dx & v=\frac{x^3}{3}
\end{matrix}\right| \to \frac{x^{3} \ln{\left(x \right)}}{3} - \frac{x^{3}}{9}+K$
\end{solution}

\part[2]
$$\int x^{2} \ln^2{\left(x \right)}\, dx$$
\begin{solution}
$\left|\begin{matrix}
u = \ln^2 x & du=\frac{2}{x}\cdot \ln x \, dx \\ dv= x^2 \, dx & v=\frac{x^3}{3}
\end{matrix}\right| \to \frac{x^{3} \ln{\left(x \right)}^{2}}{3} - \frac{2 x^{3} \ln{\left(x \right)}}{9} + \frac{2 x^{3}}{27}+K$
\end{solution}

\part[1]
$$\int\limits_{2}^{3} \frac{1}{2 x^{2} - 4 x + 2}\, dx$$
\begin{solution}
$\frac{1}{4}  \land F(x)= -\frac{1}{2 x - 2}$
\end{solution}
\end{parts}


\question[3] Calcule las dimensiones de tres campos cuadrados que no tienen ningún lado común y que
satisfacen que el perímetro de uno de ellos es triple que el de otro y, además, se necesitan 1248 metros de
valla para vallar completamente los tres campos, de manera que la suma de las áreas es la mínima posible.
\begin{solution}
    Lados= $x$, $3x$ e $y$ \\
    $4x+4(3x)+4y=1248 \to 4(x+3x+y)=1248 \to 4x+y=312 \to y=312-4x$ \\
    maximizar $f(x)=x^2+(3x)^2+(312-4x)^2=26 x^{2} - 2496 x + 97344= x^{2} - 96 x + 3744$ \\
    $f'(x)=0 \to 2x-96=0 \to x=48 \land f''(48)=2>0 \to$ en $x=48$ hay un mínimo.
    Solución: 48, 144, 120 m  
    
\end{solution}
\addpoints

\end{questions}

\end{document}
\grid
